\section{优缺点分析与结论}

\subsection{模型优点}

本文模型具有以下优点：
\begin{enumerate}
  \item 数据口径清晰，区分正向销量、净销量和负向调整。
  \item 模型与题目四问一一对应，形成门店、商品、类别和门店--商品的分层预测框架。
  \item 主模型优先采用移动平均、同星期均值、指数平滑和固定效应回归，公式简单、可解释、易复现。
  \item 对综合模型进行了信息泄露审查，并以严格 7 日递推验证作为主口径。
  \item 对外部变量和商品关联均设置了证据边界，没有把相关性写成因果性。
\end{enumerate}

\subsection{模型局限}

本文模型存在以下局限。第一，门店--商品粒度存在大量零销量与间歇性需求，76 条序列中 54 条零销量比例不低于 50\%，构成细粒度预测的结构性误差上限；WAPE 因低销量分母放大而偏高，混合策略只能部分缓解。第二，未来真实天气与促销安排不可得，含外部变量预测依赖历史同期情景，属于条件预测。第三，附件缺少顾客购物篮、价格、库存与陈列信息，无法识别商品互补/替代和促销因果效应，问题二、问题三结论均限于统计关联。第四，严格递推验证仅含 4 个完整 7 日窗口，统计功效有限，显著性检验和稳定性判断应谨慎推广。

\subsection{模型推广}

该框架可以推广到更多门店和更多品类。若新增门店或商品，只需按日期--门店--商品构造同样的日销量面板，并在门店、商品和类别层级分别保留可解释基准模型，再用严格 7 日递推检验综合模型是否真正优于强基准。若后续取得价格、库存、陈列位置和促销力度等字段，可以把这些变量并入 $x_{s,p,t}$，在预测任务中作为外部特征，在解释任务中则需继续区分统计关联和因果效应；若存在促销随机实验或明确的政策切换，再考虑因果识别方法。对零销量比例较高的商品，可进一步引入 Croston 族或 Syntetos--Boylan Approximation 等间歇需求模型 \cite{syntetos2005accuracy}，但仍应与简单指数平滑、移动平均等基准模型在同一验证口径下比较。

\subsection{结论}

本文完成了休闲零食连锁店商品销量预测的四问建模。问题一保留移动平均、同星期均值和简单指数平滑，建立了门店、商品和门店--商品的基础预测；问题二以附件类别字段整合同类零食，商品相关性分析作为同步波动证据，类别预测主方案为类别聚合后直接预测；问题三采用合并天气、对数销量和历史控制固定效应回归，得到外部变量与销量之间的条件统计关联；问题四采用综合 Ridge 模型融合前三问信息，并用严格 7 日递推验证检验未来 7 天预测能力。

本文围绕销量预测四问，建立了由简单基准到综合融合、口径逐级衔接的分层框架。最重要的发现有两点：其一，在本题数据下，简单可解释模型在严格递推口径下与复杂模型相当，复杂度的边际误差收益受限；其二，门店--商品层高误差主要由间歇需求结构决定，而非单纯建模能力不足。问题四综合 Ridge 虽然能整合前三问信息，但在严格 7 日递推主粒度下未超过问题一门店--商品简单指数平滑强基准模型。进一步的低销量鲁棒性检验表明，对低销量组合采用指数平滑、对常规组合采用综合 Ridge 的混合策略在严格递推主粒度上 WAPE=75.85\%，相对简单基准下降 0.49 个百分点，但统计检验未支持“显著改进”。基于此，本文以“低销量序列指数平滑 + 常规序列综合 Ridge”的混合策略生成未来 7 天条件预测；所有层级预测经统一整数化后由门店--商品--日期最细粒度求和，保证门店、商品、类别与总计汇总一致。整数化预测总量为 1328，汇总表见附录 D，完整明细可作为电子附件提交。
